% ---------------------------------------------------------------------------
% preface.tex
% ---------------------------------------------------------------------------
\chapter*{Preface}
\addcontentsline{toc}{chapter}{Preface}
\markboth{Preface}{Preface}

You have written VHDL for five years. You know what a signal is, why it is not
a variable, what a delta cycle does to your simulation, and how to tell
\code{numeric\_std} exactly how many bits you meant. You have argued with a
compiler that refused to add a \code{signed} to an \code{integer}, and you have
usually turned out to be the one who was wrong.

Now somebody has told you the last course of your degree is in SystemVerilog.

The bad news first. SystemVerilog is not a tidier VHDL. It comes from a
different tradition, it makes different assumptions about what a language is
for, and several of the habits that made you a careful VHDL designer will make
you a dangerous SystemVerilog one. The compiler that used to catch your
mistakes will now accept them silently. A width mismatch that VHDL would refuse
to analyse becomes a truncation that costs you an afternoon in the waveform
viewer.

The good news is larger. Almost everything you know transfers. A process is an
\code{always} block. A generic is a parameter. A package is a package. This
booklet is the translation exercise, and you will move through it quickly.

\section*{What this booklet is}

A cookbook, in the literal sense. It is organised so that you can look
something up, find the VHDL you would have written next to the SystemVerilog
you should write instead, and get on with your work. Almost every section is
built around a side-by-side comparison:

\begin{rosetta}
\begin{rleft}
process (clk) is
begin
  if rising_edge(clk) then
    q <= d;
  end if;
end process;
\end{rleft}
\begin{rright}
always_ff @(posedge clk) begin
  q <= d;
end
\end{rright}
\end{rosetta}

VHDL is always on the left, SystemVerilog always on the right. Where a
construct has no counterpart in the other language, the booklet says so
plainly rather than inventing one.

Three kinds of box interrupt the text. A \emph{trap} marks something that
specifically catches VHDL designers: not an obscure corner of the standard, but
a mistake you are statistically likely to make in your first month. A
\emph{recipe} gives a pattern to copy. A plain note supplies background you can
skip on a first reading.

\section*{What this booklet is not}

It is not a language reference. IEEE 1800-2017 is 1315 pages long and this
booklet is not a substitute for it. Where a construct has corners this text
does not explore, it says so and points you at the standard.

It is not a synthesis manual. Everything presented as synthesisable is
synthesisable with mainstream tools, but the booklet does not discuss timing
closure, physical design or vendor-specific attributes.

It is not advocacy. VHDL is an excellent language and, in several respects
catalogued honestly at the end of \cref{ch:philosophy}, a better-designed one.
The question this booklet answers is not which language is better but how to
move between them without losing the discipline the first one taught you.

It is also not the whole book. This is the \textbf{Design Cookbook}, one of
four companions covering \tool{SystemVerilog for VHDL Designers} in full.
Verification-only material (classes, randomisation, coverage, assertions,
UVM) lives in \refverificationbook; \tool{Verilator}, \tool{QuestaSim} and
Python testbenches live in \refsimulationbook; a complete worked design is
in \refexamplebook.

\section*{How this booklet is organised}

Four chapters. \Cref{ch:philosophy} explains why the two languages feel so
different, which is mostly a story about their simulation semantics and about
what each one expects a compiler to do for you; it ends with an honest account
of what VHDL does better. Read it first, even if you are tempted to skip to
the syntax. \Cref{ch:rosetta} and \cref{ch:rosetta2} are the translation
reference, split into the static half (files, ports, parameters, types,
expressions) and the behavioural half (processes, statements, generate,
instantiation, packages, time). Together they are the part of the booklet you
will keep open beside your editor, and the complete mapping table at the end
of \cref{ch:rosetta2} is the page to photocopy. \Cref{ch:beyond} covers the
design-side constructs SystemVerilog adds and VHDL has no answer to:
interfaces, packed structures, parameterised types, \code{bind}.

\section*{Acknowledgement}

The CORDIC rotator used as a running example throughout this booklet, and
developed in full in \refexamplebook, was chosen because it is small enough
to read in one sitting and large enough to need a state machine, a
parameterised datapath, a fixed-point format and a handshake protocol.
